\chapter{The Jordan Curve Theorem}\label{ch:jordan}

\begin{center}\itshape
A closed curve in the plane that does not cross itself has an inside and an outside.\\
Everyone knows it, and the proof uses nothing but continuity.
\end{center}

The closed disc is the case one can see. No continuous map of the disc onto its boundary can fix the boundary pointwise, and every continuous map of the disc to itself has a fixed point. Both facts are winding numbers. A general simple closed curve, which may have no tangent and may even have positive area, is the theorem. Its proof is the disc's theorem, applied to paths that cross a rectangle.

Camille Jordan stated the theorem in his \emph{Cours d'analyse} of 1887: a closed curve without self-intersections divides the plane into two
regions, one inside and one outside. His proof was long, and for a century it was said to be wrong; Thomas Hales, who checked a
formal proof of the theorem by computer in 2005, read it again and found it essentially correct, if incomplete for polygons, which
Jordan thought obvious. Oswald Veblen gave the first proof accepted as complete in 1905. The difficulty is that a ``curve'' may be
very wild: a continuous injective image of a circle can have positive area, as William Osgood showed in 1903, and it can have no tangent anywhere, like the boundary of Helge von Koch's
snowflake (1904).
Any proof must work with continuity alone. Volume~8 uses the inside and the outside to free Cauchy's theorem from star-shaped sets.

\section{Winding numbers of continuous curves}

Volume~5 defined the winding number $n(\gamma,a)=\frac1{2\pi i}\int_\gamma\frac{\dd z}{z-a}$ of a piecewise $C^1$ closed path. For the Jordan curve theorem we need it for
continuous curves, where there is no integral. A \emph{curve} is a continuous map $\gamma:[\alpha,\beta]\to\C$, closed if $\gamma(\alpha)=\gamma(\beta)$, with image $\gamma^*$.

\begin{lemma}[continuous argument]\label{lem:lift}
If $\gamma:[\alpha,\beta]\to\C^*$ is continuous, there is a continuous $\theta:[\alpha,\beta]\to\R$ with $\gamma(t)=|\gamma(t)|e^{i\theta(t)}$, and any two such functions differ by
a constant in $2\pi\Z$.
\end{lemma}
\begin{proof}
Let $m=\min|\gamma|>0$. By uniform continuity there is $\delta>0$ with $|\gamma(s)-\gamma(t)|<m$ when $|s-t|<\delta$, so that $\bigl|\frac{\gamma(s)}{\gamma(t)}-1\bigr|<1$: the quotient lies
in the right half-plane, where the principal logarithm $\operatorname{Log}$ is continuous. Take a subdivision $\alpha=t_0<\dots<t_N=\beta$ of mesh less than $\delta$, choose $\theta(t_0)$ with
$\gamma(t_0)=|\gamma(t_0)|e^{i\theta(t_0)}$, and define successively $\theta(t)=\theta(t_j)+\operatorname{Im}\operatorname{Log}\frac{\gamma(t)}{\gamma(t_j)}$ for $t\in[t_j,t_{j+1}]$. This is continuous and has the required
property. Two such functions differ by a continuous function with values in $2\pi\Z$, which is constant by the intermediate value theorem.
\end{proof}

\begin{definition}
For a closed curve $\gamma$ and $a\notin\gamma^*$, the \emph{winding number} is $n(\gamma,a)=\frac{\theta(\beta)-\theta(\alpha)}{2\pi}\in\Z$, where $\theta$ is a continuous argument of $\gamma-a$.
\end{definition}

For a piecewise $C^1$ path this is Volume~5's winding number: on each piece of the subdivision, $L(t)=\log|\gamma(t)-a|+i\theta(t)$ is
$\operatorname{Log}\frac{\gamma(t)-a}{\gamma(t_j)-a}$ plus a constant, so $L'=\frac{\gamma'}{\gamma-a}$, and $\int_\gamma\frac{\dd z}{z-a}=L(\beta)-L(\alpha)=i(\theta(\beta)-\theta(\alpha))$. Arguments add under
multiplication, so $n(\gamma_1\gamma_2,0)=n(\gamma_1,0)+n(\gamma_2,0)$ for closed curves in $\C^*$.

\begin{proposition}\label{prop:windingprops}
Let $\gamma$ and $\gamma_0,\gamma_1$ be closed curves on $[\alpha,\beta]$, and $a\in\C$.
\begin{enumerate}
\item[(a)] (\emph{Dog on a leash.}) If $|\gamma_1(t)-\gamma_0(t)|<|\gamma_0(t)-a|$ for all $t$, then $n(\gamma_1,a)=n(\gamma_0,a)$.
\item[(b)] (\emph{Homotopy.}) If $H:[0,1]\times[\alpha,\beta]\to\C\setminus\{a\}$ is continuous and each $H(s,\cdot)$ is closed, then $n(H(s,\cdot),a)$ does not depend on $s$.
\item[(c)] $a\mapsto n(\gamma,a)$ is constant on each component of $\C\setminus\gamma^*$, and $0$ on the unbounded one.
\end{enumerate}
\end{proposition}
\begin{proof}
(a) The quotient $q=\frac{\gamma_1-a}{\gamma_0-a}=1+\frac{\gamma_1-\gamma_0}{\gamma_0-a}$ lies in the disc $D(1,1)$, so $\operatorname{Im}\operatorname{Log}q$ is a continuous argument of $q$, and it takes the same
value at $\alpha$ and $\beta$. So $n(q,0)=0$, and $n(\gamma_1,a)=n(q,0)+n(\gamma_0,a)$.

(b) $H$ is uniformly continuous on the compact square and $m=\min|H-a|>0$, so there is $\delta$ with $|H(s,t)-H(s',t)|<m$ for $|s-s'|<\delta$; by (a), $s\mapsto n(H(s,\cdot),a)$ is
locally constant on $[0,1]$, hence constant.

(c) If $|a'-a|<d(a,\gamma^*)$, then $|(\gamma-a')-(\gamma-a)|<|\gamma-a|$, and (a) gives $n(\gamma,a')=n(\gamma,a)$. A locally constant function is constant on connected sets. If
$|a|>\max|\gamma|$, then $\operatorname{Re}\bigl((\gamma(t)-a)\overline a\bigr)\le|\gamma(t)||a|-|a|^2<0$: the curve $-\frac{\overline a}{|a|}(\gamma-a)$ stays in the right half-plane, where
$\operatorname{Im}\operatorname{Log}$ is a continuous argument taking the same value at both ends. So $n(\gamma,a)=0$, and $n(\gamma,\cdot)$ vanishes on the unbounded component, which contains such points.
\end{proof}

\begin{theorem}[no retraction; Brouwer]\label{thm:brouwer}
\begin{enumerate}
\item[(a)] There is no continuous map $r:\overline D\to\partial D$ of the closed unit disc onto its boundary with $r(z)=z$ on $\partial D$.
\item[(b)] Every continuous map of a closed disc, or of a closed square, into itself has a fixed point.
\end{enumerate}
\end{theorem}
\begin{proof}
(a) $H(s,t)=r(se^{it})$, $0\le s\le1$, $0\le t\le2\pi$, is a homotopy in $\C\setminus\{0\}$ from the constant curve $r(0)$, of winding number $0$, to $e^{it}$, of winding number $1$.

(b) Suppose $f:\overline D\to\overline D$ has no fixed point. For $z\in\overline D$ let $u=z-f(z)\ne0$ and let $r(z)=z+\tau u$, where $\tau\ge0$ is the larger root of $|z+\tau u|^2=1$:
\[\tau(z)=\frac{-\operatorname{Re}(\overline zu)+\sqrt{(\operatorname{Re}\overline zu)^2+|u|^2(1-|z|^2)}}{|u|^2}.\]
This is continuous, and $r(z)$ is where the ray from $f(z)$ through $z$ leaves the disc. If $|z|=1$, then $\operatorname{Re}(\overline zu)=1-\operatorname{Re}(\overline zf(z))\ge0$, so $\tau=0$ and $r(z)=z$. This
contradicts (a). A square is mapped onto a disc by the homeomorphism $z\mapsto\frac{\|z\|_\infty}{|z|}z$ (and $0\mapsto0$), which carries the problem over.
\end{proof}

The next lemma, due to Ryuji Maehara (1984), is the bridge between Brouwer and Jordan. It says that a path crossing a rectangle from left to right
and a path crossing it from top to bottom must meet, the topological content of the game of Hex, which can never end in a draw.

\begin{lemma}[crossing paths]\label{lem:crossing}
Let $R=[x_0,x_1]\times[y_0,y_1]$, and let $h,v:[-1,1]\to R$ be continuous with $h_1(-1)=x_0$, $h_1(1)=x_1$, $v_2(-1)=y_0$, $v_2(1)=y_1$ (indices denote coordinates).
In words: a path $h$ from the left side of $R$ to the right side and a path $v$ from the bottom to the top must meet. Then $h$ and $v$ meet: $h(s)=v(t)$ for some $s,t$.
\end{lemma}
\begin{proof}
If not, the function $N(s,t)=\|h(s)-v(t)\|_\infty$ is positive everywhere, and
\[F(s,t)=\Bigl(\frac{v_1(t)-h_1(s)}{N(s,t)},\frac{h_2(s)-v_2(t)}{N(s,t)}\Bigr)\]
is a continuous map of the square $[-1,1]^2$ into its boundary. By Theorem~\ref{thm:brouwer}(b) it has a fixed point $(s,t)$, and $|s|=1$ or $|t|=1$. If $s=-1$, then
$F_1=\frac{v_1(t)-x_0}N\ge0$, while $F_1=s=-1$; if $s=1$, then $F_1=\frac{v_1(t)-x_1}N\le0\ne1$; the cases $t=\pm1$ are the same with the second coordinate. Contradiction.
\end{proof}

\section{The Jordan curve theorem}

A \emph{Jordan curve} is the image $J$ of a closed curve $\gamma:[\alpha,\beta]\to\C$ that is injective on $[\alpha,\beta[$; equivalently, a subset of the plane homeomorphic to a circle. An
\emph{arc} is a subset homeomorphic to $[0,1]$. Since $J$ is compact, $\C\setminus J$ is open; its connected components are open and path-connected, exactly one of them is
unbounded (it contains the complement of a large disc, which is connected), and the boundary of each component lies in $J$.

\begin{lemma}[arcs do not separate]\label{lem:arc}
The complement of an arc is connected.
\end{lemma}
\begin{proof}
Let $A$ be an arc and $\phi:A\to[0,1]$ a homeomorphism. Suppose that $\C\setminus A$ is not connected. One of its components is unbounded, so another one, $W$, is bounded. The boundary of
$W$ lies in $A$: a point $z\in\partial W$ outside $A$ would lie in some component $W'$ of $\C\setminus A$; this open set meets $W$, so $W'=W$ and $z\in W$, impossible for a boundary point of an open
set.

By Tietze's extension theorem (Exercise~\ref{ex:tietze}, applied to $2\phi-1$), $\phi$ extends to a continuous $\Phi:\C\to[0,1]$, and $\rho=\phi^{-1}\circ\Phi$ is a continuous map $\C\to A$ that is
the identity on $A$. Choose $z_1\in W$ and a closed disc $\Delta=\overline D(z_1,R)$ whose interior contains the compact set $A\cup\overline W$. Define $g:\Delta\to\C$ by
\[
g(z)=\rho(z)\ \text{ for }z\in\overline W,\qquad g(z)=z\ \text{ for }z\in\Delta\setminus W.
\]
The two closed sets $\overline W$ and $\Delta\setminus W$ cover $\Delta$ and meet in $\partial W\subset A$, where $\rho(z)=z$; so $g$ is well defined and continuous. It never takes the value $z_1$:
on $\overline W$ its values lie in $A$, which does not meet $W$, and on $\Delta\setminus W$ it is the identity. Let $\pi(z)=z_1+R\frac{z-z_1}{|z-z_1|}$ be the radial projection onto the circle
$\partial\Delta$. Then $\pi\circ g$ is a continuous map of $\Delta$ onto $\partial\Delta$, equal to the identity on $\partial\Delta$ (which lies in $\Delta\setminus W$). After the similarity
$z\mapsto\frac{z-z_1}R$, this contradicts Theorem~\ref{thm:brouwer}(a).
\end{proof}

\begin{theorem}[Jordan curve theorem]\label{thm:jordan}
Let $J$ be a Jordan curve. Then $\C\setminus J$ has exactly two components, one bounded and one unbounded, and $J$ is the boundary of each of them.
\end{theorem}

The bounded component is the \emph{inside} of $J$, the unbounded one the \emph{outside}. The proof is Maehara's. It builds, out of $J$ and three vertical segments, a path that cuts a
rectangle from top to bottom, and then uses Lemma~\ref{lem:crossing} three times: any path from the left to the right of the rectangle that avoids it must be impossible. Figure~\ref{fig:jordan}
shows the construction.

\begin{figure}[ht]
\centering
\fig{ch09-jordan}
\caption{Maehara's construction. The arc $J_n$ (charcoal) contains the highest point $l$ of $J$ on the segment $[n,s]$; $m$ is its lowest point there. The arc $J_s$ (ochre) first meets $[m,s]$ at
$p$, and $q$ is its lowest point on $[n,s]$. The top-to-bottom path $\Gamma$ (coral) runs down from $n$ to $l$, along $J_n$ to $m$, across to $p$, along $J_s$ to $q$ and down to $s$. The
midpoint $z_0$ of $[m,p]$ is inside.}
\label{fig:jordan}
\end{figure}

\begin{proof}
\emph{Setting up.} Since $J$ is compact, there are $a,b\in J$ at maximal distance. A similarity of the plane maps Jordan curves to Jordan curves and preserves everything in the statement, so
we may assume $a=-1$ and $b=1$. Every $z\in J$ then satisfies $|z-a|\le2$ and $|z-b|\le2$. So $J$ lies in the rectangle $R=[-1,1]\times[-2,2]$; it meets the boundary of $R$ only at $a$ and $b$, because
$|\operatorname{Im}z|\le\sqrt3<2$ on the lens $\{|z-a|\le2,\ |z-b|\le2\}$, and a point of the lens with $\operatorname{Re}z=-1$ satisfies $|z-b|\ge2$, with equality only at $a$ (likewise for $b$). The points $a$, $b$
cut $J$ into two arcs from $a$ to $b$ that meet only at $a$ and $b$. Let $n=2i$ and $s=-2i$ be the midpoints of the top and bottom sides of $R$. Each arc goes from the left side of $R$ to the right
side, and $[n,s]$ from the top to the bottom, so by Lemma~\ref{lem:crossing} each arc meets $[n,s]$. ``Higher'' means a larger imaginary part.

\emph{Step 1: a point $z_0$.} Let $l$ be the highest point of the compact set $J\cap[n,s]$; let $J_n$ be the arc containing $l$, and $J_s$ the other one ($l\ne a,b$, so $l$ lies on only one arc). Let $m$
be the lowest point of $J_n\cap[n,s]$. The segment $[m,s]$ meets $J_s$: otherwise the path that goes down $[n,l]$, along $J_n$ from $l$ to $m$, and down $[m,s]$ would cross $R$ from top to bottom
without meeting $J_s$. (The segment $[n,l]$ meets $J$ only at $l$, which is not on $J_s$; the part of $J_n$ between $l$ and $m$ does not contain $a$ or $b$.) This contradicts Lemma~\ref{lem:crossing}
for $J_s$. So let $p$ be the highest point of $J_s\cap[m,s]$ and $q$ the lowest point of $J_s\cap[n,s]$. The open segment $]m,p[$ misses $J_n$ (by the choice of $m$) and $J_s$ (by the choice of
$p$). Let $z_0$ be its midpoint, and $U$ the component of $\C\setminus J$ containing $z_0$.

\emph{Step 2: $U$ is bounded.} If not, there is a path in $U$ from $z_0$, an interior point of $R$, to a point outside $R$. Let $w$ be the first point of $\partial R$ on it, and $\delta$ the part of
the path from $z_0$ to $w$: it lies in $R$ and misses $J$, and $w\ne a,b$. The set $\partial R\setminus\{a,b\}$ has two pieces, the upper one $\partial^+$ through $n$ and the lower one $\partial^-$
through $s$, both disjoint from $J$. If $w\in\partial^-$, go down $[n,l]$, along $J_n$ to $m$, along $[m,z_0]$, along $\delta$ to $w$, and along $\partial^-$ to $s$: a top-to-bottom path in $R$ that
misses $J_s$, which contradicts Lemma~\ref{lem:crossing}. If $w\in\partial^+$, go along $\partial^+$ from $n$ to $w$, back along $\delta$ to $z_0$, along $[z_0,p]$, along $J_s$ from $p$ to $q$, and down
$[q,s]$: a top-to-bottom path that misses $J_n$ (note that $[q,s]$ lies below $m$), again a contradiction. So $U$ is bounded, and $\C\setminus J$ has at least two components.

\emph{Step 3: $J$ is the boundary of every component.} Let $W$ be a component; as in the proof of Lemma~\ref{lem:arc}, $\partial W\subset J$. If $\partial W\ne J$, choose $x\in J\setminus\partial W$.
Since $\partial W$ is closed, a small open subarc of $J$ around $x$ misses it, so $\partial W$ lies in the arc $A$ obtained by removing that subarc from $J$. By Lemma~\ref{lem:arc}, the open set
$\C\setminus A$ is connected, hence path-connected (Proposition~\ref{prop:openconnected}). Take $w\in W$ and $w'$ in another component (Step 2), and a path $c:[0,1]\to\C\setminus A$ from $w$ to $w'$. Let
$t^*=\sup\{t:c([0,t])\subset W\}$. Then $c(t^*)\in\overline W$, and $c(t^*)\notin W$, since otherwise the path would stay in the open set $W$ a little longer (and $t^*<1$ because $w'\notin W$). So
$c(t^*)\in\partial W\subset A$, a contradiction.

\emph{Step 4: $U$ is the only bounded component.} Let $\Gamma$ be the path that goes down $[n,l]$, along $J_n$ to $m$, along $[m,p]$, along $J_s$ to $q$, and down $[q,s]$. Its points outside $J$
lie on $[n,l[$, on $]m,p[$ and on $]q,s]$. The middle segment lies in $U$. The first one is connected, misses $J$ and contains $n$, which is joined to infinity by the upward vertical ray, which misses
$J$; so it lies in the unbounded component, and so does $]q,s]$. Now let $W$ be a bounded component other than $U$. Then $\Gamma$ does not meet $W$. Neither $a$ nor $b$ lies on the compact
set $\Gamma$, so some discs $D(a,\varepsilon)$ and $D(b,\varepsilon)$ miss $\Gamma$; by Step 3, $a,b\in\partial W$, so there are $a'\in W\cap D(a,\varepsilon)$ and $b'\in W\cap D(b,\varepsilon)$. The set $\C\setminus R$ is
connected, unbounded and disjoint from $J$, so it lies in the unbounded component, and $W\subset R$. The path $[a,a']$, then a path in $W$ from $a'$ to $b'$, then $[b',b]$, goes from the left of
$R$ to the right inside $R$ and misses $\Gamma$, which crosses $R$ from top to bottom. This contradicts Lemma~\ref{lem:crossing} once more.
\end{proof}

\begin{think}
Where did the proof use that $J$ has no self-intersections? Draw a figure eight and follow the four steps: which one breaks down? And why could we not simply say ``the winding number is $1$
inside and $0$ outside''?
\end{think}

\begin{remark}
The inside $U$ of a Jordan curve can be strange. Osgood's curves of positive area have insides whose boundary carries area, and in 1910 Brouwer found three disjoint open sets in the plane
with the same boundary, the ``Lakes of Wada'', so the Jordan curve theorem really does use that $J$ is a curve. Schoenflies proved in 1906 that $U$ is nevertheless homeomorphic to a disc,
by a homeomorphism of the whole plane carrying $J$ to the unit circle. The conformal version of this, Carath\'eodory's theorem, needs the Riemann mapping theorem, which Volume~9 proves, and a boundary analysis, which is not given here.
\end{remark}


% Polygonal Jordan theorem.
\section{Jordan's theorem for polygons}

A \emph{polygon} $P$ is a closed path $v_1v_2\cdots v_nv_1$ made of $n\ge3$ segments (the \emph{edges} $e_i=[v_i,v_{i+1}]$, indices modulo $n$). It is \emph{simple} if two edges meet only when
they are consecutive, and then only at their common vertex; consecutive edges are not collinear. We identify $P$ with the union of its edges, a compact subset of $\R^2$.

\begin{theorem}[Jordan's theorem for polygons]\label{thm:jordanpoly}
If $P$ is a simple polygon, then $\R^2\setminus P$ has exactly two connected components. One is bounded and the other is not, and $P$ is the boundary of each.
\end{theorem}

\begin{figure}[ht]\centering\fig{ch10-polygon}
\caption{The parity of the number of crossings of a ray with $P$ tells inside from outside: $3$ for the point $x$, $2$ for the point $y$.}\label{fig:polygon}
\end{figure}

The proof has two halves. Counting crossings of a ray (Figure~\ref{fig:polygon}) shows that there are \emph{at least} two components; walking along both sides of $P$ shows that there are \emph{at most} two.
Here ``connected'' means path-connected, and the components of an open subset of the plane are open.

\subsection*{At least two components: a parity}

Choose a unit vector $d$ that is parallel to no edge and to no line through two vertices; only finitely many directions are excluded. For a vertex $v$, call the line $v+\R d$ a \emph{bad line}; there are $n$ of them,
and each contains exactly one vertex. For $x\notin P$ off the bad lines, the ray $x+\R_{\ge0}d$ meets no vertex; it meets each edge in at most one point (it is not parallel to the edge), and there it crosses the
edge from one side to the other. Let $c(x)$ be the number of edges it meets, and $\pi(x)=c(x)\bmod2$.

\begin{lemma}\label{lem:parity}
Every $x\in\R^2\setminus P$ has a disc neighbourhood $D$, disjoint from $P$, on which $\pi$ is constant (on the points of $D$ off the bad lines).
\end{lemma}

\begin{proof}
Let $x\notin P$. The ray from $x$ meets at most one vertex. Let $\rho>0$ be smaller than the distance from $x$ to $P$, and from the ray at $x$ to every vertex not on it. Then for $x'$ in the disc $D$ of
radius $\rho$ about $x$, the ray from $x'$ is a translate of the ray from $x$ by less than $\rho$.

If the ray from $x$ meets no vertex, then (taking $\rho$ smaller if needed, by compactness of the edges and transversality) the ray from $x'$ meets exactly the edges met by the ray from $x$, and $\pi(x')=\pi(x)$. If it meets the vertex $v$,
between the edges $e$ and $f$, the same holds for all edges except $e$ and $f$. Points $x'$ of $D$ off the bad line through $v$ lie on one of its two sides, and the ray from $x'$ passes $v$ on that side. An edge at $v$ is met by
the ray from $x'$ exactly when the edge leaves $v$ towards the side of the bad line where $x'$ is... more simply: when $x'$ crosses the bad line, the ray sweeps across $v$, and it stops meeting the edges at $v$ on one side and starts
meeting those on the other side. If $e$ and $f$ lie on the same side of the bad line, the count of these two edges changes from $0$ to $2$ or from $2$ to $0$; if they lie on opposite sides, it stays $1$. In all cases $\pi$ does not change.
\end{proof}

By the lemma, $\pi$ extends to a function on all of $\R^2\setminus P$ that is constant on a neighbourhood of every point: put $\pi(x)=\pi(x')$ for any $x'$ near $x$ off the bad lines. A locally constant function is constant along
any path (the set of times where it equals its initial value is open and closed in $[0,1]$), hence on each component.

Now take a point $m$ in the interior of an edge $e$, off the bad lines, and $\varepsilon>0$ so small that the segment $[m-\varepsilon d,m+\varepsilon d]$ meets $P$ only at $m$ and its endpoints are off the bad lines. The ray from $m-\varepsilon d$ is the
segment $[m-\varepsilon d,m+\varepsilon d)$ followed by the ray from $m+\varepsilon d$, and it meets $P$ in one more point, $m$. So $\pi(m-\varepsilon d)\ne\pi(m+\varepsilon d)$: these two points lie in different components, and there are at least two.

\subsection*{At most two components: the two sides}

Let $r>0$ be smaller than half the length of every edge and than half the distance from each vertex to every edge not containing it (positive, since $P$ is simple). Let $B_v$ be the open disc of radius $r$ about the vertex $v$.
In polar coordinates about $v$, $B_v\cap P$ is two radii, in the directions of the two edges at $v$. They cut $B_v\setminus P$ into two open circular sectors, each path-connected (it is an interval of angles times $(0,r)$).

For an edge $e=[v,w]$, let $e'$ be the part of $e$ at distance at least $\frac r2$ from both $v$ and $w$. It is a compact segment disjoint from all other edges, so it has positive distance $\delta_e$ from them. Let $R_e$ be the open rectangle of points whose orthogonal projection on the line of $e$ lies in the interior of $e'$ and whose distance to that line is less than $h$, with $h<\min(\delta_e,\frac r4)$ and so small that the ends of $R_e$ near $v$ and $w$ stay inside $B_v$ and $B_w$ away from the other edge there. Then $R_e\cap P=e'$, and $R_e\setminus P$ is two open
rectangles $R_e^+$ and $R_e^-$, one on each side of $e$, each convex.

\begin{lemma}\label{lem:sides}
Each of $R_e^\pm$ meets exactly one sector at each end; at each vertex, the two sides of an edge meet different sectors.
\end{lemma}

\begin{proof}
The end of $R_e^+$ near $v$ is a connected subset of $B_v\setminus P$, so it lies in one sector. In polar coordinates about $v$, the points of $R_e^\pm$ near $v$ have angles just above and just below the direction of $e$. The two directions of the edges at $v$ cut the circle of
angles into two arcs, and the direction of $e$ is an endpoint of both; angles just above it lie in one arc and angles just below in the other. So $R_e^+$ and $R_e^-$ meet different sectors.
\end{proof}

Build a graph whose $4n$ nodes are the $2n$ sides $R_e^\pm$ and the $2n$ sectors, with an edge when they intersect. By Lemma~\ref{lem:sides}, each side has degree $2$ (one sector at each end), and each sector, at a vertex where the edges $e$ and $f$ meet, contains exactly one side of $e$ and one side of $f$, so it has degree $2$ too.
A graph in which every node has degree $2$ is a union of disjoint cycles. Start from $R_{e_1}^+$ and walk: sector at $v_2$, a side of $e_2$, sector at $v_3$, and so on around $P$, back to a side of $e_1$ after visiting one side of every edge. Either the walk returns to $R_{e_1}^+$, and the
other $2n$ nodes form a second cycle, or it returns to $R_{e_1}^-$ and continues through the remaining nodes. So there are at most two cycles. The union of the nodes of a cycle is path-connected (each node is, and consecutive nodes intersect), so we get at most two path-connected sets $S_1,S_2\subseteq\R^2\setminus P$ (possibly equal).

The sets $B_v$ and $R_e$ are open and cover $P$, and each of them minus $P$ is contained in $S_1\cup S_2$. Let $U$ be a component of $\R^2\setminus P$. It is open and not the whole plane, and the plane is connected, so $U$ has a boundary point $p$, which lies in $P$ (a point of another
component has a neighbourhood disjoint from $U$). The open set of the cover containing $p$ meets $U$, so $U$ meets $S_1$ or $S_2$, and therefore contains it. Hence there are at most two components.

\begin{proof}[Proof of Theorem~\ref{thm:jordanpoly}]
There are at least two components and at most two, so exactly two, and each contains one of $S_1$, $S_2$; in particular $S_1\ne S_2$, so we are in the case of two cycles, and each component contains one side of every edge and one sector at every vertex. Hence every point of $P$ is a limit
of points of each component, and the boundary of each component, which lies in $P$, is all of $P$. Finally, $P$ lies in a disc $B$; the complement of $B$ is path-connected and disjoint from $P$, so it lies in one component, the unbounded one, and the other component is contained in $B$.
\end{proof}

\begin{remark}[the general theorem]\label{rem:jordangeneral}
For an arbitrary continuous simple closed curve the theorem is true but much harder: one approximates the curve by polygons, and the difficulty is to control the components in the limit. Hales's 2007 formal proof in HOL Light follows this route,
and his article on Jordan's original text explains which steps Jordan took for granted. The polygonal case already shows what readers of Jordan argued about: the treatment of rays through vertices (Lemma~\ref{lem:parity}), which Jordan regarded as obvious.
\end{remark}

Look back at the proof of Lemma~\ref{lem:parity}: it contains a sentence that starts one way and then says ``more simply''. That is how first drafts look, and it is the first thing to fix (Exercise~\ref{ex:rewritelemma}).

\begin{think}
Where exactly was simplicity of $P$ used? Try the proof on a figure eight and see which step fails, and how many components there are.
\end{think}

\begin{beauty}
Jordan's theorem is the archetype of a statement that is obvious and hard. The polygonal proof is honest about why: all the difficulty sits at the vertices, in one parity check and one
picture of two sectors, and once those are written down carefully there is nothing left to believe. Writing well is the art of putting the difficulty where the reader can see it.
\end{beauty}

\begin{exercise}[$\star$]\label{ex:rewritelemma}
Rewrite the second paragraph of the proof of Lemma~\ref{lem:parity} so that it says what it means the first time, with the case analysis at the vertex stated as a separate claim.
\end{exercise}


\section*{Hints}

\begin{hint}{ex:rewritelemma}
Isolate the claim: ``as $x'$ crosses the bad line through $v$, the number of edges at $v$ met by the ray changes by $0$ or $\pm2$''.
\end{hint}

\section*{Solutions}

\begin{solution}{ex:rewritelemma}
\idea State the local claim first, then use it.
\writeup ``Suppose the ray from $x$ meets the vertex $v$, where the edges $e$ and $f$ meet. Every other edge is met by the ray from $x'$ if and only if it is met by the ray from $x$. \emph{Claim:} the number of edges among $e,f$ met by the ray from $x'$ is the same, modulo $2$, on both sides of the bad line through $v$. Indeed, the ray from $x'$ is parallel to the bad line and passes $v$ on the side of $x'$, at small distance; so an edge at $v$ is met by it exactly when the edge lies on the side of $x'$. If $e$ and $f$ lie on the same side, the count is $0$ on one side of the bad line and $2$ on the other; if they lie on opposite sides, it is $1$ on both. Hence $\pi$ is constant on $D$.'' The original sentence tried to give the reason and the conclusion at once; separating the claim from its proof is almost always the fix.
\end{solution}


\section{Exercises}

\begin{exercise}[$\star$]\label{ex:ftawind}
\emph{The fundamental theorem of algebra, by winding.} Let $p(z)=z^n+c_{n-1}z^{n-1}+\dots+c_0$ with $n\ge1$, and $\gamma_R(t)=Re^{it}$, $0\le t\le2\pi$. (a) Show that $n(p\circ\gamma_R,0)=n$ when $R>1+\sum_j|c_j|$.
(b) If $p$ had no zero, what would $H(s,t)=p(sRe^{it})$ say? Conclude.
\end{exercise}

\section{Hints}

\begin{hint}{ex:ftawind}
On a big circle, $p$ is a dog on a short leash held by $z^n$. Then shrink the circle to its centre.
\end{hint}

\section{Solutions}

\begin{solution}{ex:ftawind}
\idea Winding numbers are robust: a curve close to $z^n$ winds $n$ times, and a curve that can be shrunk to a point without passing through $0$ winds $0$ times.
\writeup (a) On $|z|=R$ with $R>1$, $|p(z)-z^n|\le\sum_j|c_j|R^j\le R^{n-1}\sum_j|c_j|<R^n=|z^n|$. By Proposition~\ref{prop:windingprops}(a), $n(p\circ\gamma_R,0)=n(\gamma_R^n,0)=n$, since $nt$ is a
continuous argument of $R^ne^{int}$. (b) If $p$ never vanishes, $H$ is a continuous map $[0,1]\times[0,2\pi]\to\C\setminus\{0\}$ through closed curves, from the constant curve $p(0)$, of winding number $0$, to
$p\circ\gamma_R$. By Proposition~\ref{prop:windingprops}(b), $0=n$, which is false. So $p$ has a zero.
\end{solution}

\section*{Sources and further reading}
\addcontentsline{toc}{section}{Sources and further reading}

C.~Jordan, \emph{Cours d'analyse de l'\'Ecole polytechnique}, vol.~3 (Gauthier-Villars, 1887); O.~Veblen, ``Theory of plane curves in non-metrical analysis situs'', \emph{Trans.\ Amer.\ Math.\ Soc.}
6 (1905), 83--98; W.~F.~Osgood, ``A Jordan curve of positive area'', \emph{Trans.\ Amer.\ Math.\ Soc.} 4 (1903), 107--112. T.~C.~Hales, ``The Jordan curve theorem, formally and informally'', \emph{Amer.\ Math.\
Monthly} 114 (2007), 882--894, tells the history and the computer verification, and his ``Jordan's proof of the Jordan curve theorem'', \emph{Studies in Logic, Grammar and Rhetoric}
10(23) (2007), 45--60, rehabilitates Jordan. H.~von Koch's snowflake is in \emph{Arkiv f\"or Mat., Astron.\ och Fys.} 1 (1904), 681--702. The Lakes of Wada are in L.~E.~J.~Brouwer, ``Zur Analysis
Situs'', \emph{Math.\ Ann.} 68 (1910), 422--434; for Schoenflies's theorem see C.~Thomassen, ``The Jordan--Sch\"onflies theorem and the classification of surfaces'', \emph{Amer.\ Math.\ Monthly} 99
(1992), 116--130. The proof given here is R.~Maehara, ``The Jordan curve theorem via the Brouwer fixed point theorem'', \emph{Amer.\
Math.\ Monthly} 91 (1984), 641--643.
